% GENERATED FILE. Edit canonical lessons, then run npm run generate:books.
% canonical-source-hash: fnv1a64:22624141730c9aac
% canonical-lessons: HI-C429-parci, HI-C429-kard, HI-C429-dastkhat, HI-C429-mudra, HI-C429-bhunana, HI-R429-review-422-429

\chapter{Slips, Cards and Signatures}
\label{ch:hi-slips-cards-and-signatures}
\hlchaptermodality{429}

\begin{chapteropening}
\textbf{By the end of this chapter:} \emph{I can fill in a slip, pay by card, sign my name and cash a cheque.}

Complete the last lesson of chapter 429: I can fill in a slip, pay by card, sign my name and cash a cheque.
\end{chapteropening}

\section[\texorpdfstring{\hi{पर्ची} (parcī)}{पर्ची (parcī)}]{\hi{पर्ची} (parcī) — a slip (of paper), a chit}

\label{lesson:HI-C429-parci}



\begin{quote}
Before the new one: say the Hindi for a budget, then the Hindi for to collect (money owed).
\end{quote}

\subsection*{You'll want to know: \hi{पर्ची}}
\textbf{\hi{पर्ची}} — \emph{parcī} — \textquotedblleft{}a slip (of paper), a chit\textquotedblright{}.

\begin{quote}
\hi{बैंक} \hi{में} \hi{पहले} \hi{एक} \hi{पर्ची} \hi{भरिए।}
\end{quote}

\begin{quote}
\emph{baiṅk meṁ pahle ek parcī bhariye.}
\end{quote}

\begin{quote}
At the bank, fill in a slip first.
\end{quote}

Feminine.

\begin{cousinweb}
\textbf{\hi{पर्ची}} (\emph{parcī}) is usually traced to Persian \emph{parča}, \textquotedblleft{}a piece, a scrap\textquotedblright{}, with the Hindi ending -ī for a small one: a slip for a bank, a doctor or a lottery.
\end{cousinweb}

\begin{grammarlens}[title={What you've built}]
One more word for getting around.
\end{grammarlens}

\subsection*{Your turn}
\begin{itemize}
\raggedright
  \item \emph{Say it:} \emph{parcī}
  \item \emph{Say it:} \emph{parcī}, once more
  \item \emph{Say it:} \emph{baiṅk meṁ pahle ek parcī bhariye.} — At the bank, fill in a slip first.
  \item \emph{Recall:} say the Hindi for a budget, then the Hindi for to collect (money owed), then say \emph{parcī} again
\end{itemize}

\begin{tcolorbox}[breakable,colback=teal!4,colframe=teal!35!black,title={Before you move on}]
What does \hi{पर्ची} mean? (\textquotedblleft{}A slip (of paper), a chit\textquotedblright{}.) Say it once more.
\end{tcolorbox}
\section[\texorpdfstring{\hi{कार्ड} (kārḍ)}{कार्ड (kārḍ)}]{\hi{कार्ड} (kārḍ) — a card}

\label{lesson:HI-C429-kard}



\begin{quote}
Before the new one: say the Hindi for to collect (money owed), then the Hindi for a slip.
\end{quote}

\subsection*{You'll want to know: \hi{कार्ड}}
\textbf{\hi{कार्ड}} — \emph{kārḍ} — \textquotedblleft{}a card\textquotedblright{}.

\begin{quote}
\hi{क्या} \hi{आप} \hi{कार्ड} \hi{से} \hi{पैसे} \hi{लेते} \hi{हैं}?
\end{quote}

\begin{quote}
\emph{kyā āp kārḍ se paise lete haiṁ?}
\end{quote}

\begin{quote}
Do you take cards?
\end{quote}

Masculine.

\begin{cousinweb}
\textbf{\hi{कार्ड}} (\emph{kārḍ}) is borrowed from English \emph{card}: a bank card, a wedding card, a visiting card.
\end{cousinweb}

\begin{grammarlens}[title={What you've built}]
One more word for getting around.
\end{grammarlens}

\subsection*{Your turn}
\begin{itemize}
\raggedright
  \item \emph{Say it:} \emph{kārḍ}
  \item \emph{Say it:} \emph{kārḍ}, once more
  \item \emph{Say it:} \emph{kyā āp kārḍ se paise lete haiṁ?} — Do you take cards?
  \item \emph{Recall:} say the Hindi for to collect (money owed), then the Hindi for a slip, then say \emph{kārḍ} again
\end{itemize}

\begin{tcolorbox}[breakable,colback=teal!4,colframe=teal!35!black,title={Before you move on}]
What does \hi{कार्ड} mean? (\textquotedblleft{}A card\textquotedblright{}.) Say it once more.
\end{tcolorbox}
\section[\texorpdfstring{\hi{दस्तख़त} (dastkhat)}{दस्तख़त (dastkhat)}]{\hi{दस्तख़त} (dastkhat) — a signature (everyday)}

\label{lesson:HI-C429-dastkhat}



\begin{quote}
Before the new one: say the Hindi for a slip, then the Hindi for a card.
\end{quote}

\subsection*{You'll want to know: \hi{दस्तख़त}}
\textbf{\hi{दस्तख़त}} — \emph{dastkhat} — \textquotedblleft{}a signature (everyday)\textquotedblright{}.

\begin{quote}
\hi{यहाँ} \hi{अपने} \hi{दस्तख़त} \hi{कीजिए।}
\end{quote}

\begin{quote}
\emph{yahā̃ apne dastkhat kījiye.}
\end{quote}

\begin{quote}
Please sign here.
\end{quote}

Masculine.

Usually plural in agreement: \hi{मैंने} \hi{दस्तख़त} \hi{किए}, I signed.

\begin{cousinweb}
\textbf{\hi{दस्तख़त}} (\emph{dastkhat}) is usually traced to Persian \emph{dastkhaṭṭ}, from \emph{dast}, \textquotedblleft{}hand\textquotedblright{}, and \emph{khaṭṭ}, \textquotedblleft{}writing\textquotedblright{}: your hand's own mark. \hi{हस्ताक्षर} is the formal word.
\end{cousinweb}

\begin{grammarlens}[title={What you've built}]
One more word for getting around.
\end{grammarlens}

\subsection*{Your turn}
\begin{itemize}
\raggedright
  \item \emph{Say it:} \emph{dastkhat}
  \item \emph{Say it:} \emph{dastkhat}, once more
  \item \emph{Say it:} \emph{yahā̃ apne dastkhat kījiye.} — Please sign here.
  \item \emph{Recall:} say the Hindi for a slip, then the Hindi for a card, then say \emph{dastkhat} again
\end{itemize}

\begin{tcolorbox}[breakable,colback=teal!4,colframe=teal!35!black,title={Before you move on}]
What does \hi{दस्तख़त} mean? (\textquotedblleft{}A signature (everyday)\textquotedblright{}.) Say it once more.
\end{tcolorbox}
\section[\texorpdfstring{\hi{मुद्रा} (mudrā)}{मुद्रा (mudrā)}]{\hi{मुद्रा} (mudrā) — a currency}

\label{lesson:HI-C429-mudra}



\begin{quote}
Before the new one: say the Hindi for a card, then the Hindi for a signature (everyday).
\end{quote}

\subsection*{You'll want to know: \hi{मुद्रा}}
\textbf{\hi{मुद्रा}} — \emph{mudrā} — \textquotedblleft{}a currency\textquotedblright{}.

\begin{quote}
\hi{हर} \hi{देश} \hi{की} \hi{अपनी} \hi{मुद्रा} \hi{होती} \hi{है।}
\end{quote}

\begin{quote}
\emph{har deś kī apnī mudrā hotī hai.}
\end{quote}

\begin{quote}
Every country has its own currency.
\end{quote}

Feminine.

\begin{cousinweb}
\textbf{\hi{मुद्रा}} (\emph{mudrā}) is a tatsama, Sanskrit \emph{mudrā}, \textquotedblleft{}a seal, a stamp\textquotedblright{} (Monier-Williams); a stamped coin led to money in general, and modern Hindi uses it for a currency.
\end{cousinweb}

\begin{grammarlens}[title={What you've built}]
One more word for getting around.
\end{grammarlens}

\subsection*{Your turn}
\begin{itemize}
\raggedright
  \item \emph{Say it:} \emph{mudrā}
  \item \emph{Say it:} \emph{mudrā}, once more
  \item \emph{Say it:} \emph{har deś kī apnī mudrā hotī hai.} — Every country has its own currency.
  \item \emph{Recall:} say the Hindi for a card, then the Hindi for a signature (everyday), then say \emph{mudrā} again
\end{itemize}

\begin{tcolorbox}[breakable,colback=teal!4,colframe=teal!35!black,title={Before you move on}]
What does \hi{मुद्रा} mean? (\textquotedblleft{}A currency\textquotedblright{}.) Say it once more.
\end{tcolorbox}
\section[\texorpdfstring{\hi{भुनाना} (bhunānā)}{भुनाना (bhunānā)}]{\hi{भुनाना} (bhunānā) — to cash; to change (a note)}

\label{lesson:HI-C429-bhunana}



\begin{quote}
Before the new one: say the Hindi for a signature (everyday), then the Hindi for a currency.
\end{quote}

\subsection*{You'll want to know: \hi{भुनाना}}
\textbf{\hi{भुनाना}} — \emph{bhunānā} — \textquotedblleft{}to cash; to change (a note)\textquotedblright{}.

\begin{quote}
\hi{मैंने} \hi{बैंक} \hi{में} \hi{चेक} \hi{भुनाया।}
\end{quote}

\begin{quote}
\emph{maiṁne baiṅk meṁ cek bhunāyā.}
\end{quote}

\begin{quote}
I cashed the cheque at the bank.
\end{quote}

Past, with \hi{ने} (he or she …): \hi{उसने} \hi{सौ} \hi{का} \hi{नोट} \hi{भुनाया}, \emph{usne sau kā noṭ bhunāyā}.

\begin{cousinweb}
\textbf{\hi{भुनाना}} (\emph{bhunānā}) is listed by CDIAL under Sanskrit \emph{bhajyate}, \textquotedblleft{}is broken\textquotedblright{} (CDIAL 9361), with Hindi \emph{bhunnā}, \textquotedblleft{}to be changed (of money)\textquotedblright{}, and its causative \emph{bhunānā}: to break a big note into smaller money. CDIAL calls the u of these forms unexplained.
\end{cousinweb}

\begin{grammarlens}[title={What you've built}]
One more word for getting around.
\end{grammarlens}

\subsection*{Your turn}
\begin{itemize}
\raggedright
  \item \emph{Say it:} \emph{bhunānā}
  \item \emph{Say it:} \emph{bhunānā}, once more
  \item \emph{Say it:} \emph{maiṁne baiṅk meṁ cek bhunāyā.} — I cashed the cheque at the bank.
  \item \emph{Recall:} say the Hindi for a signature (everyday), then the Hindi for a currency, then say \emph{bhunānā} again
\end{itemize}

\begin{tcolorbox}[breakable,colback=teal!4,colframe=teal!35!black,title={Before you move on}]
What does \hi{भुनाना} mean? (\textquotedblleft{}To cash; to change (a note)\textquotedblright{}.) Say it once more.
\end{tcolorbox}
\section[(dialogue)]{Looking back: from Houses Large and Small to Slips, Cards and Signatures}

\label{lesson:HI-R429-review-422-429}



\begin{quote}
Say the Hindi for a hut and a bungalow.
\end{quote}

\subsection*{Your turn}
Say these aloud:

\begin{itemize}
\raggedright
  \item a hut, a bungalow, a house (the building), a compound, a thatched roof
  \item a guard, a pole, noise, to growl, to bark
  \item a roadside eatery, a sweet-maker, a handcart, a hawker, a betel leaf
  \item an uproar, to chase away, a sewer, a stink, to seep
  \item a sum, in cash, a cheque, a balance, a cash machine
  \item credit, an instalment, an account, a broker, to pay off
  \item to earn, to blow, a profit (in trade), a budget, to collect (money owed)
  \item a slip, a card, a signature (everyday), a currency, to cash
\end{itemize}

\begin{tcolorbox}[breakable,colback=teal!4,colframe=teal!35!black,title={Before you move on}]
A hut? (\textbf{\hi{झोंपड़ी}}, \emph{jhõpṛī}.) A bungalow? (\textbf{\hi{कोठी}}, \emph{koṭhī}.) A house (the building)? (\textbf{\hi{मकान}}, \emph{makān}.) A compound? (\textbf{\hi{अहाता}}, \emph{ahātā}.) A thatched roof? (\textbf{\hi{छप्पर}}, \emph{chappar}.) A guard? (\textbf{\hi{गार्ड}}, \emph{gārḍ}.) A pole? (\textbf{\hi{खंभा}}, \emph{khambhā}.) Noise? (\textbf{\hi{शोर}}, \emph{śor}.) To growl? (\textbf{\hi{गुर्राना}}, \emph{gurrānā}.) To bark? (\textbf{\hi{भौंकना}}, \emph{bhaũknā}.) A roadside eatery? (\textbf{\hi{ढाबा}}, \emph{ḍhābā}.) A sweet-maker? (\textbf{\hi{हलवाई}}, \emph{halvāī}.) A handcart? (\textbf{\hi{ठेला}}, \emph{ṭhelā}.) A hawker? (\textbf{\hi{फेरीवाला}}, \emph{pherīvālā}.) A betel leaf? (\textbf{\hi{पान}}, \emph{pān}.) An uproar? (\textbf{\hi{हल्ला}}, \emph{hallā}.) To chase away? (\textbf{\hi{खदेड़ना}}, \emph{khadeṛnā}.) A sewer? (\textbf{\hi{गटर}}, \emph{gaṭar}.) A stink? (\textbf{\hi{बदबू}}, \emph{badbū}.) To seep? (\textbf{\hi{रिसना}}, \emph{risnā}.) A sum? (\textbf{\hi{रक़म}}, \emph{raqam}.) In cash? (\textbf{\hi{नक़द}}, \emph{naqad}.) A cheque? (\textbf{\hi{चेक}}, \emph{cek}.) A balance? (\textbf{\hi{बैलेंस}}, \emph{bailens}.) A cash machine? (\textbf{\hi{एटीएम}}, \emph{eṭīem}.) Credit? (\textbf{\hi{उधार}}, \emph{udhār}.) An instalment? (\textbf{\hi{क़िस्त}}, \emph{qist}.) An account? (\textbf{\hi{हिसाब}}, \emph{hisāb}.) A broker? (\textbf{\hi{दलाल}}, \emph{dalāl}.) To pay off? (\textbf{\hi{चुकाना}}, \emph{cukānā}.) To earn? (\textbf{\hi{कमाना}}, \emph{kamānā}.) To blow? (\textbf{\hi{फूँकना}}, \emph{phū̃knā}.) A profit (in trade)? (\textbf{\hi{मुनाफ़ा}}, \emph{munāfā}.) A budget? (\textbf{\hi{बजट}}, \emph{bajaṭ}.) To collect (money owed)? (\textbf{\hi{वसूलना}}, \emph{vasūlnā}.) A slip? (\textbf{\hi{पर्ची}}, \emph{parcī}.) A card? (\textbf{\hi{कार्ड}}, \emph{kārḍ}.) A signature (everyday)? (\textbf{\hi{दस्तख़त}}, \emph{dastkhat}.) A currency? (\textbf{\hi{मुद्रा}}, \emph{mudrā}.) To cash? (\textbf{\hi{भुनाना}}, \emph{bhunānā}.)
\end{tcolorbox}
